\chapter{Responsible Use of Resources and End-of-Life Strategy }
\section*{Responsible Use of Resources}
The hydrogel bioprinter was developed with the goal of using efficient materials and reusable components. The majority of the mechanical system consists of standard parts, including aluminium extrusions, TR8x8 lead screws, GT2 belts and pulleys, NEMA17 stepper motors, and linear bearings, which can all be reused in future student projects. 6 small PLA parts were 3D printed (4 feet and 2 limit switch mounts), resulting in minimal filament waste. The printer also includes several custom aluminium components (such as the syringe and bed plate) to ensure durability, most of which are made from aluminium plates to minimize machine waste.

The custom PCB was designed to maximise modularity, re-usability, and upgradability, by incorporating 5 stepper ports with breakout driver headers (StepStick form factor), 12 V switching ports, EEPROM storage, USB and SD card interface, and various breakout pins ensuring that the circuit board can be adapted across many different embedded systems applications. The experimental phase used store-bought agar powder as a hydrogel surrogate gel that is safe, biodegradable, inexpensive, and did not require a special disposal method.

\section*{End-of-Life Strategy}

At the end of the bioprinter's operational life, it can be disassembled into its mechanical and electronic components. The aluminium extrusions, motors, and other standard parts can be reused in future university projects or recycled through standard metal recycling facilities (along with the rest of the custom aluminium components). The PCB can be re-purposed to operate in other mechatronic systems or responsibly disposed of as electronic-waste and the PLA components can be recycled. The bioprinter’s modular design ensures that its components can be repaired, upgraded, or replaced, thereby extending the system's operational life and minimising environmental impact.